\subsection{\bi{全状态比较与端点取值}{Global comparison and evaluation at the endpoint}}
\bi{为说明第二个不等式怎样从逐步比较产生，记 $K_cF(r)=\mathbb EF(r-Z_c)$，$c=-\log r_{\rm action}$，$v_t=K_a^t\phi_n$，$f_t$ 为同一固定动作和的密度。取 $b=m(a)-a$。记 $J_t$ 为参考核的 Poisson 校正，$s_t$ 为可求和的误差系数。内侧比较函数为}{To derive the second inequality from one-step comparisons write $K_cF(r)=\mathbb EF(r-Z_c)$, with $c$ the negative logarithm of the action, $v_t=K_a^t\phi_n$, and $f_t$ the fixed-action sum density. Set $b=m(a)-a$. With $J_t$ the reference-kernel Poisson correction and $s_t$ a bounded accumulated error coefficient, the inner comparison function is}
\[
 I_t=v_t\exp\{(\alpha J_t+s_t f_t)/v_t\}+\Lambda tb^2/n .
\]
\bi{外侧取两支单步上解的最小值 $O_t$。以 $\eta_\mu(r/t)$ 记非负鞍点参数（$r/t\le\mu$ 时置零）。在 $\eta_\mu\le2d(\alpha-1)$ 上取 $W_t=\min(O_t,I_t)$，在其余处取 $W_t=O_t$。附录中自然管与正中偏差带的核估计、早期左带与晚期左带的比较，给出下列拼接：}{Let $O_t$ be the minimum of two outer supersolutions and let $\eta_\mu(r/t)$ denote the nonnegative saddle parameter, taken to be zero when $r/t\le\mu$. Use $W_t=\min(O_t,I_t)$ for $\eta_\mu\le2d(\alpha-1)$ and $W_t=O_t$ elsewhere. The tube and right moderate-deviation estimates, together with the early and late left-band comparisons in the supplement, give the following gluing argument.}

\bi{在正侧接口 $d(\alpha-1)\le\eta_\mu\le2d(\alpha-1)$ 已有 $I_t\ge O_t$，因此选外支。接口以下的当前点所有前态均满足 $r-Z_c\le r$；从 $t$ 到 $t-1$ 的鞍点增量为 $O(1/t)$，故这些前态均允许内支。自然管 $|r-tm(a)|\le A\sqrt{t\log t}$ 由内侧全动作比较控制，右侧剩余点使用正中偏差比较。左侧在 $t\le\sqrt n$ 时，将 $r\le tm(a)-C t^{2/3}$ 划给外支，相邻带使用负倾斜比较。在 $t\ge n^{1/3}$ 时，$r\le tm(a)-A'\sqrt{t\log n}$ 的 Chernoff 尾被 $tb^2/n$ 项吸收，余下左侧处在加宽的自然管内。两个时间范围重叠。常数在 $n$ 前固定；这些局部比较对所有动作成立，故}{On the positive interface $d(\alpha-1)\le\eta_\mu\le2d(\alpha-1)$ one has $I_t\ge O_t$, so the outer branch applies. Below the interface every predecessor satisfies $r-Z_c\le r$, and changing $t$ to $t-1$ changes the saddle parameter by $O(1/t)$ on the relevant compact range. The inner branch is therefore available at every predecessor. Use the all-action inner estimate in $|r-tm(a)|\le A\sqrt{t\log t}$ and the positive moderate-deviation estimate on its right. For $t\le\sqrt n$, the far left $r\le tm(a)-Ct^{2/3}$ uses the outer branch, with a negative-tilt comparison on the adjoining band. For $t\ge n^{1/3}$ a Chernoff tail on $r\le tm(a)-A'\sqrt{t\log n}$ is absorbed by $tb^2/n$; the rest lies in an enlarged natural tube. These time ranges overlap. All constants are fixed before $n$. The resulting one-step comparisons hold for every action:}
\[
 \sup_{c\ge0}K_c W_{t-1}\le W_t,\qquad t_0<t\le n .
\]
\bi{初值使用整个函数的平移。由单步密度界，任意策略 $k$ 步和的尾至多为 $\Gamma(k,1)$ 尾的 $2^k$ 倍。Gamma 风险率递增，且 $\Pr(\Gamma(k,1)>4k)\le(4e^{-3})^k<2^{-k}$，故该和依随机序受 $\Gamma(k,1)+4k$ 控制。参考动作的和支配 $\Gamma(k,1)$。由于 $\phi_n$ 递减且 $W_{t_0}\ge v_{t_0}$，固定平移 $h_0=4t_0$ 即控制初值。平移与所有 $K_c$ 交换，归纳给 $\mathcal U_n(s)\le W_n(s-h_0)\le I_n(s-h_0)$，末个比较只在临界窗口使用。}{Initialize by translating the whole function. The increment density bound gives a $k$-step policy tail at most $2^k$ times the $\Gamma(k,1)$ tail. The gamma hazard is increasing and $\Pr(\Gamma(k,1)>4k)\le(4e^{-3})^k<2^{-k}$, so the policy sum is stochastically dominated by $\Gamma(k,1)+4k$. The reference sum dominates $\Gamma(k,1)$. Since $\phi_n$ decreases and $W_{t_0}\ge v_{t_0}$, the fixed translation $h_0=4t_0$ controls the initial value. Translation commutes with every $K_c$. Induction gives $\mathcal U_n(s)\le W_n(s-h_0)\le I_n(s-h_0)$, using the last comparison only in the critical window.}

\bi{匹配点处 Poisson 校正的二次部分为零；余项为 $f_n(s)(p_a/8)(\mu+o(1))(\sum_{j=1}^n1/j-1)$，其中 $p_a=(1+e^a)^{-1}\to(\mu-1)/\mu$。故 $J_n(s)=(g+o(1))f_n(s)L$。固定平移的代价为 $O(f_n(s))$，$b^2=O(L/n)$，$s_n=O(1)$。由于 $f_n(s)L\to0$ 且 $v_n(s)\to1$，指数展开的二次项为 $o(f_n(s)L)$，给出所需上界。附录的证明次序及精确 Poisson 恒等式详见随稿技术材料。}{At the matching point the quadratic part of the Poisson correction vanishes. The remaining term is $f_n(s)(p_a/8)(\mu+o(1))(\sum_{j=1}^n1/j-1)$ with $p_a=(1+e^a)^{-1}\to(\mu-1)/\mu$, giving $J_n(s)=(g+o(1))f_n(s)L$. The fixed translation costs $O(f_n(s))$, while $b^2=O(L/n)$ and $s_n=O(1)$. Since $f_n(s)L\to0$ and $v_n(s)\to1$, the quadratic remainder in the exponential is $o(f_n(s)L)$, as required. The accompanying technical supplement gives the underlying local estimates and the exact Poisson identity.}
